Every number of the manuscript is produced by a script of the repository from the committed artifacts of the case C04 (the documents under \texttt{data/derived/C04}) or from an experiment of the study, and no raw CEREP file is needed to rebuild any of them. The engine is \texttt{riskvalidation} in release 1.1.0 with the product's generators and risk-weight function, and the code is that of the repository named in the header of the document, in the release that carries this document version; the seeds are those stated in Section~\ref{sec:design-protocol} and in the experiment sources. From the manuscript folder, with the project's Python environment (each script finds the repository root from its own location), five experiment scripts write their outputs as compact JSON in the folder \texttt{experiments}:
\begin{itemize}
\item \texttt{python experiments/exp\_definitions.py} (the real-data definitions, the bootstrap and the Fitch years with a populated default category, under a minute);
\item \texttt{python experiments/exp\_withdrawals.py} (the reproduction of the committed rungs and the extension, about \minWd{} minutes);
\item \texttt{python experiments/exp\_theory.py} (the two-state check, about \minTheory{} minute);
\item \texttt{python experiments/exp\_thin.py} (the enumeration, under a minute);
\item \texttt{python experiments/exp\_extra.py} (the duration estimator and the three further transition tests, on the seeds of the extension, about \minExtra{} minutes).
\end{itemize}
\begin{sloppypar}
Two generators then write the rest: \texttt{python figures/make\_figures.py} writes the eight vector figures from those outputs and from the committed documents, and \texttt{python tables/make\_tables.py} writes every table and the file \texttt{numbers.tex}, whose macros hold the numbers quoted in the prose; its twin \texttt{numbers.json} gives the source of each macro. The document is built with \texttt{pdflatex}, \texttt{bibtex} and two further passes of \texttt{pdflatex} from the manuscript folder.
\end{sloppypar}
